%==============================================================================
% PARTIDA 01: CONEXION A LA RED EXTERNA DE MEDIDORES
% Codigo: OE.5.1.1.1
%==============================================================================
\subsection{OE.5.1.1.1 Acometida e Instalacion de Caja Portamedidor Monofasico}

\subsubsection{Especificaciones Tecnicas}

La acometida comprende el conjunto de elementos comprendidos entre la red publica de distribucion y el medidor de energia, incluyendo el cableado de alimentacion, la caja portamedidor monofasica y sus accesorios. La instalacion debe cumplir con lo establecido en el CNE-U (Seccion 2: Reglas para instalaciones de acometida) y la Norma EM.010 del RNE.

La caja portamedidor debe ser del tipo monofasico, con tapa precintable, de material termoplastico o metalico resistente a la intemperie, con grado de proteccion IP-54 minimo. Debera alojar el medidor de energia y permitir las conexiones de entrada y salida de los conductores.

El conductor de acometida sera de cobre tipo LSOH (Low Smoke Zero Halogen) de 10 mm2 para fase y neutro, con aislamiento para 0.6/1 kV, con capacidad de corriente suficiente para soportar la maxima demanda del proyecto (21.73 A), incluyendo un factor de seguridad y reserva para ampliaciones futuras.

\subsubsection{Requisitos de Materiales}

\begin{itemize}
  \item Caja portamedidor monofasica termoplastica o metalica, con tapa precintable y visor.
  \item Conductor de cobre LSOH 10 mm2 (fase y neutro), color negro y blanco respectivamente, para 0.6/1 kV.
  \item Conductor de cobre desnudo 6 mm2 para conexion a tierra del sistema de medicion.
  \item Tuberia PVC pesada (PV-P) de 1'' de diametro para proteccion del alimentador desde la acometida hasta el medidor.
  \item Conectores de compression tipo bimetálico para conexiones en la caja portamedidor.
  \item Cinta aislante autofundente y cinta vinilica para sellado de conexiones.
  \item Abrazaderas y soportes tipo "U" galvanizados para fijacion de tuberia a muros.
  \item Pernos de anclaje y tarugos de expansion para fijacion de la caja.
  \item Gel siliconico o grasa dielectrico para proteccion de conexiones contra humedad.
\end{itemize}

\subsubsection{Metodo de Ejecucion}

\begin{enumerate}
  \item Realizar el replanteo y marcar la ubicacion de la caja portamedidor en la fachada o limite de propiedad, segun indicaciones de la empresa distribuidora (Electro Puno S.A.A.) y el plano IE-01.
  \item Fijar la caja portamedidor al muro mediante pernos de anclaje, asegurando que quede nivelada y firmemente sujeta. La altura recomendada es entre 1.50 m y 1.80 m sobre el nivel del piso.
  \item Instalar la tuberia PVC pesada (PV-P) de 1'' desde el punto de acometida hasta la caja portamedidor, asegurando una pendiente minima del 0.5\% para evitar acumulacion de humedad.
  \item Pasar los conductores de fase y neutro de 10 mm2 a traves de la tuberia, dejando una longitud adicional de 0.50 m en cada extremo para conexiones.
  \item Conectar los conductores en la caja portamedidor utilizando conectores de compression, asegurando un apriete adecuado y aplicando gel dielectrico para proteccion contra corrosion.
  \item Instalar el conductor de puesta a tierra de 6 mm2 desde la barra de tierra del tablero general hasta el electrodo de puesta a tierra, pasando por la caja portamedidor.
  \item Sellar las entradas de la caja portamedidor con cinta aislante y compuesto sellador para evitar el ingreso de humedad y polvo.
  \item Rotular la caja portamedidor con la identificacion del suministro y el numero de medidor asignado.
\end{enumerate}

\subsubsection{Control de Calidad}

\begin{itemize}
  \item Verificar que la caja portamedidor cuente con certificacion de producto y cumpla con las normas NTP correspondientes.
  \item Comprobar la continuidad electrica de los conductores de fase, neutro y tierra antes de conectar el medidor.
  \item Medir la resistencia de aislamiento entre conductores fase-neutro y fase-tierra, debiendo ser superior a 1 MOhm con megohmetro de 1000 V.
  \item Inspeccionar visualmente que no existan empalmes intermedios en los conductores dentro de la tuberia.
  \item Verificar la correcta fijacion y nivelacion de la caja portamedidor.
  \item Comprobar que los precintos de seguridad esten correctamente instalados.
  \item Realizar prueba de funcionamiento del medidor una vez energizado el suministro, verificando la lectura correcta de consumo.
\end{itemize}

\subsubsection{Metodo de Medicion}

La medicion se realizara por unidad (UND) de caja portamedidor monofasico completamente instalada, incluyendo todos los accesorios, conexiones, sellados y pruebas. Se medira por cada punto de suministro instalado y verificado.

\subsubsection{Unidad de Medida}

UND (Unidad) -- Cada caja portamedidor monofasico instalada y operativa.

\subsubsection{Forma de Pago}

El pago se efectuara por unidad completamente instalada, aprobada por la supervision y verificada mediante las pruebas de calidad correspondientes. Incluye el suministro de materiales, mano de obra, herramientas, equipos y todos los recursos necesarios para dejar la partida en condiciones de operacion.
